\documentclass[11pt]{article}
\usepackage[margin=2.5cm]{geometry}
\usepackage{amsmath}
\usepackage{amssymb}
\usepackage{booktabs}
\usepackage{hyperref}

% Numbers are computed in partB.ipynb and cross-checked in
% ../code/html_checks/tangency_check.py and borrowing_check.py.
% No local LaTeX installation: compile on Overleaf.

\newcommand{\rborrow}{r_{\text{borrow}}}

\title{FINA 3103 -- Homework 1, Part B}
\author{Erik Jin (21434993) \and Yining Wang (21279642) \and Ziyou Chen (21435052)}
\date{October 2026}

\begin{document}
\maketitle

\section*{Question 1. Your own risk aversion and the two-step workflow}

\subsection*{Inputs}

The certainty equivalents (CEs) our group submitted in Part~A, used exactly as submitted:

\begin{center}
\begin{tabular}{llccc}
\toprule
Bet & Outcomes & Erik & Elaine & Chen \\
\midrule
X & 50\%: $+30\%$, 50\%: $-10\%$ & 10\% & 4\% & 7\% \\
Y & 80\%: $+15\%$, 20\%: $-25\%$ & 12\% & 3\% & 5\% \\
Z & 20\%: $+39\%$, 80\%: $-1\%$  & 5\%  & 3.5\% & 2\% \\
\bottomrule
\end{tabular}
\end{center}

Elaine is Yining Wang; Erik is Erik Jin and Chen is Ziyou Chen.

The capital market consists of a stock fund S ($E[r_S] = 10\%$, $\sigma_S = 18\%$), a bond fund B ($E[r_B] = 6\%$, $\sigma_B = 10\%$), correlation $\rho_{SB} = 0.25$ and a risk-free rate $r_f = 3\%$.

%------------------------------------------------------------------------------
\subsection*{(a) Risk aversion from the certainty equivalents}

For a bet with outcomes $r_i$ and probabilities $p_i$,
\[
E[r] = \sum_i p_i r_i, \qquad \sigma = \sqrt{\sum_i p_i \bigl(r_i - E[r]\bigr)^2}.
\]
Under mean-variance utility the CE equals the utility of the bet, $CE = E[r] - \tfrac12 A\sigma^2$, which gives
\[
A = \frac{2\bigl(E[r] - CE\bigr)}{\sigma^2}.
\]

\begin{center}
\begin{tabular}{lcc}
\toprule
Bet & $E[r]$ & $\sigma$ \\
\midrule
X & 10\% & 20\% \\
Y & 7\%  & 16\% \\
Z & 7\%  & 16\% \\
\bottomrule
\end{tabular}
\qquad
\begin{tabular}{lccc}
\toprule
$A$ & Bet X & Bet Y & Bet Z \\
\midrule
Erik   & 0.00 & $-3.91$ & 1.56 \\
Elaine & 3.00 & 3.13    & 2.73 \\
Chen   & 1.50 & 1.56    & 3.91 \\
\bottomrule
\end{tabular}
\end{center}

\paragraph{What should hold for Bets Y and Z.}
Bets Y and Z have the same expected return (7\%) and the same standard deviation (16\%). Under mean-variance utility an investor only cares about these two moments, so each person should state the \emph{same} CE for Y and Z, and the implied $A$ should be the same for all three bets.

\paragraph{Does it hold for our group?}
No. All three members give different CEs for Y and Z:
\begin{itemize}
  \item Erik: 12\% for Y and 5\% for Z. His CE for Y is above its expected return, which implies $A_Y = -3.91$ (risk-loving), while $A_X = 0$ (risk-neutral) and $A_Z = 1.56$ (risk-averse).
  \item Elaine: 3\% for Y and 3.5\% for Z. Her implied $A$ is fairly stable (2.73--3.13), so mean-variance is a reasonable approximation for her.
  \item Chen: 5\% for Y and 2\% for Z, so his $A$ more than doubles from Y to Z.
\end{itemize}

\paragraph{Which assumption this contradicts.}
Y and Z differ only in their \emph{skewness}: Y has a small probability of a large loss (negative skew), Z a small probability of a large gain (positive skew). Since our CEs differ between them, we care about more than mean and variance. Our answers therefore contradict the assumption that returns are (approximately) normally distributed, or equivalently that investors care only about the first two moments. Elaine values the positively skewed Z more, which is the usual preference for skewness, while Erik and Chen prefer the negatively skewed Y. The variation of $A$ across bets also contradicts the assumption of a single, constant risk-aversion coefficient per investor. Elaine's answers are the exception: her implied $A$ stays between 2.73 and 3.13, so her answers broadly support the mean-variance assumption, whereas Erik's and Chen's clearly contradict it.
% TODO: match the wording to the list on the slide "What can alter or undo this result?"

%------------------------------------------------------------------------------
\subsection*{(b) Optimal risky portfolio and complete portfolios}

\paragraph{The risky portfolio $P^*$.}
All investors choose the risky portfolio with the steepest capital allocation line, i.e.\ the highest Sharpe ratio. With excess returns $R_S = E[r_S] - r_f$, $R_B = E[r_B] - r_f$ and $\sigma_{SB} = \rho_{SB}\sigma_S\sigma_B$, the first-order condition of $\max (E[r_P] - r_f)/\sigma_P$ gives $\mathbf w^* \propto \Sigma^{-1}\boldsymbol\mu$, which for two assets becomes
\[
w_S^* = \frac{R_S\,\sigma_B^2 - R_B\,\sigma_{SB}}{R_S\,\sigma_B^2 + R_B\,\sigma_S^2 - (R_S + R_B)\,\sigma_{SB}}, \qquad w_B^* = 1 - w_S^*.
\]
With $R_S = 0.07$, $R_B = 0.03$, $\sigma_S^2 = 0.0324$, $\sigma_B^2 = 0.01$ and $\sigma_{SB} = 0.0045$:
\[
w_S^* = \frac{0.000565}{0.001222} = 0.4624, \qquad w_B^* = 0.5376,
\]
\[
E[r_{P^*}] = 7.85\%, \qquad \sigma_{P^*} = 10.98\%, \qquad \text{Sharpe ratio} = \frac{7.85\% - 3\%}{10.98\%} = 0.4417.
\]

\paragraph{Capital allocation.}
The complete portfolio C puts a fraction $y$ in $P^*$ and $1-y$ in the risk-free asset, so $E[r_C] = r_f + y(E[r_{P^*}] - r_f)$ and $\sigma_C = y\sigma_{P^*}$. Maximising $U(y) = E[r_C] - \tfrac12 A\sigma_C^2$ gives the first-order condition
\[
\frac{dU}{dy} = \bigl(E[r_{P^*}] - r_f\bigr) - A y\,\sigma_{P^*}^2 = 0
\quad\Longrightarrow\quad
y^* = \frac{E[r_{P^*}] - r_f}{A\,\sigma_{P^*}^2}.
\]
The complete-portfolio weights are $(y^* w_S^*,\; y^* w_B^*,\; 1 - y^*)$. Using each member's $A$ from Bet X:

\begin{center}
\begin{tabular}{lcccccccc}
\toprule
Member & $A_X$ & $y^*$ & $w_S$ & $w_B$ & $w_{r_f}$ & $w_S/w_B$ & $E[r_C]$ & $\sigma_C$ \\
\midrule
Erik   & 0.00 & $\infty$ & $\infty$ & $\infty$ & $-\infty$ & 0.8600 & -- & -- \\
Elaine & 3.00 & 1.3410 & 0.6200 & 0.7210 & $-0.3410$ & 0.8600 & 9.50\% & 14.72\% \\
Chen   & 1.50 & 2.6820 & 1.2401 & 1.4420 & $-1.6820$ & 0.8600 & 16.01\% & 29.45\% \\
\bottomrule
\end{tabular}
\end{center}

Erik's CE for Bet X equals its expected return, so $A_X = 0$: he is risk-neutral and $y^*$ is unbounded. He would borrow without limit at $r_f$ to buy $P^*$. Elaine and Chen both have $y^* > 1$ and finance the extra position by borrowing at the risk-free rate.

\paragraph{Same ratio of S to B.}
The column $w_S/w_B$ is $0.8600 = 0.4624/0.5376$ for every member: risk aversion only changes how much is invested in $P^*$, never the mix of S and B inside it. This illustrates the \emph{separation property} (two-fund separation): the choice of the optimal risky portfolio is independent of preferences, and the investor's risk aversion only determines the capital allocation between $P^*$ and the risk-free asset.

%------------------------------------------------------------------------------
\subsection*{(c) Borrowing at 6\%, lending at 3\%}

\paragraph{Lending versus borrowing.}
The sign of the risk-free weight $1-y$ determines which side of the risk-free asset the investor is on:
\begin{itemize}
  \item \textbf{Lending} ($y \le 1$, $w_{r_f} = 1-y \ge 0$): the investor puts only part of her wealth in $P^*$ and invests the rest in the risk-free asset (e.g.\ T-bills or a deposit). She \emph{earns} $r_f = 3\%$ on that part.
  \item \textbf{Borrowing} ($y > 1$, $w_{r_f} = 1-y < 0$): the investor wants more than 100\% of her wealth in $P^*$ and finances the excess $y-1$ with a loan (e.g.\ buying on margin). She \emph{pays} $\rborrow = 6\%$ on the borrowed amount.
\end{itemize}
In (b) both sides used the same rate, so borrowing was simply a negative position in the risk-free asset and the CAL was one straight line. Now the rate depends on the sign of $1-y$. The extra cost applies only to the borrowed amount, not to the whole portfolio: with $y = 1.5$, for example, the investor pays $0.5 \times 6\% = 3\%$ of her wealth in interest instead of $0.5 \times 3\% = 1.5\%$. Investors who lend are unaffected. The expected return of the complete portfolio therefore becomes
\[
E[r_C] =
\begin{cases}
r_f + y\bigl(E[r_{P^*}] - r_f\bigr), & y \le 1, \\[2pt]
\rborrow + y\bigl(E[r_{P^*}] - \rborrow\bigr), & y > 1,
\end{cases}
\qquad \sigma_C = y\,\sigma_{P^*}.
\]
The CAL has a kink at $P^*$: its slope falls from $0.4417$ to $(7.85\% - 6\%)/10.98\% = 0.1684$. Applying the first-order condition on each segment gives
\[
y_L = \frac{E[r_{P^*}] - r_f}{A\sigma_{P^*}^2}, \qquad
y_B = \frac{E[r_{P^*}] - \rborrow}{A\sigma_{P^*}^2},
\qquad
y^* =
\begin{cases}
y_L & \text{if } y_L \le 1 \text{ (lends, unaffected)}, \\
1   & \text{if } y_B < 1 < y_L \text{ (at the kink)}, \\
y_B & \text{if } y_B \ge 1 \text{ (borrows at 6\%)}.
\end{cases}
\]

\paragraph{Who is affected.}
All three members borrowed in (b), so all three are affected.

\begin{center}
\begin{tabular}{lccccccccc}
\toprule
Member & $A_X$ & $y_L$ & $y_B$ & $y^*$ & $w_S$ & $w_B$ & $w_{r_f}$ & $U$ before & $U$ after \\
\midrule
Erik   & 0.00 & $\infty$ & $\infty$ & $\infty$ & $\infty$ & $\infty$ & $-\infty$ & $\infty$ & $\infty$ \\
Elaine & 3.00 & 1.3410 & 0.5114 & 1.0000 & 0.4624 & 0.5376 & 0.0000 & 0.0625 & 0.0604 \\
Chen   & 1.50 & 2.6820 & 1.0228 & 1.0228 & 0.4729 & 0.5499 & $-0.0228$ & 0.0950 & 0.0695 \\
\bottomrule
\end{tabular}
\end{center}

\begin{itemize}
  \item \textbf{Elaine} would borrow at 3\% ($y_L = 1.34$) but not at 6\% ($y_B = 0.51$). She ends up exactly at the kink, 100\% in $P^*$, with $E[r_C] = 7.85\%$, $\sigma_C = 10.98\%$ and utility $0.0604$ (down from $0.0625$).
  \item \textbf{Chen} still borrows, but only $2.28\%$ of his wealth instead of $168\%$. His new portfolio has $E[r_C] = 7.89\%$, $\sigma_C = 11.23\%$ and utility $0.0695$ (down from $0.0950$).
  \item \textbf{Erik} is risk-neutral. Since $E[r_{P^*}] = 7.85\% > 6\%$, borrowing is still profitable in expectation and his optimal leverage remains unbounded.
\end{itemize}

\paragraph{Range of $A$ for holding exactly 100\% in $P^*$.}
The investor neither borrows nor lends when $y_B \le 1 \le y_L$:
\[
\frac{E[r_{P^*}] - \rborrow}{\sigma_{P^*}^2} \le A \le \frac{E[r_{P^*}] - r_f}{\sigma_{P^*}^2}
\quad\Longrightarrow\quad
1.5343 \le A \le 4.0230.
\]
Elaine ($A = 3$) lies inside this range; Chen ($A = 1.5$) lies just below it and therefore still borrows a little.

%------------------------------------------------------------------------------
\subsection*{(d) Should more risk-averse investors hold more bonds?}

\paragraph{Consistency with (b).}
No. In (b) every member holds S and B in the same ratio, $w_S/w_B = 0.86$, and risk aversion only changes $y^*$. Under mean-variance utility with a single borrowing and lending rate, a more risk-averse investor should hold more \emph{cash}, not more bonds. The advisers' rule contradicts the separation property.

\paragraph{Reason 1: Borrowing is more expensive than lending (our numbers from (c)).}
Separation requires that investors can borrow and lend at the same rate. With $\rborrow = 6\%$, leverage on $P^*$ only earns $1.85\%$ per unit of borrowed money. An investor who wants more risk than $P^*$ can instead obtain it without borrowing, by tilting the risky portfolio towards stocks. Maximising $E[r] - \tfrac12 A\sigma^2$ over the stock-bond mix with $y = 1$ gives
\[
w_S = \frac{(E[r_S] - E[r_B])/A + \sigma_B^2 - \sigma_{SB}}{\sigma_S^2 + \sigma_B^2 - 2\sigma_{SB}},
\]
which is optimal for $0.7243 \le A \le 4.0230$ (between $P^*$ and the tangency portfolio relative to $6\%$):

\begin{center}
\begin{tabular}{lcccccc}
\toprule
& \multicolumn{3}{c}{$P^*$ fixed, (c)} & \multicolumn{3}{c}{Mix chosen freely, $y = 1$} \\
\cmidrule(lr){2-4}\cmidrule(lr){5-7}
Member & $w_S$ & $w_B$ & $U$ & $w_S$ & $w_B$ & $U$ \\
\midrule
Elaine ($A = 3$)   & 0.4624 & 0.5376 & 0.0604 & 0.5639 & 0.4361 & 0.0609 \\
Chen ($A = 1.5$)   & 0.4729 & 0.5499 & 0.0695 & 0.9631 & 0.0369 & 0.0757 \\
\bottomrule
\end{tabular}
\end{center}

Both are better off leaving $P^*$, and the optimal mix now depends on $A$: Chen, the less risk-averse, holds almost only stocks, while Elaine holds a larger share of bonds. With a borrowing spread the advisers' rule is therefore rational.

\paragraph{Reason 2: Preferences depend on more than mean and variance (our numbers from (a)).}
Bets Y and Z have identical mean and variance, yet every member gives them different CEs (Erik 12\% vs 5\%, Elaine 3\% vs 3.5\%, Chen 5\% vs 2\%). We care about skewness and tail risk, and $P^*$ is optimal only for an investor who looks at the first two moments. An investor who dislikes left-tail risk, as Elaine does by valuing Z above Y, would rationally hold less of the asset with larger crash risk (typically stocks) and more bonds than in $P^*$.

\paragraph{Other reasons.}
An investor with future liabilities or a long horizon may hold long-term bonds as a hedge, since they are close to risk-free for that horizon; labour income that is correlated with the stock market makes stocks riskier for that investor; and estimation error in expected returns makes the mean-variance weights unreliable in practice.

\end{document}
